\section{Experimental Details}

\subsection{Data}\label{sec:data}
For unsupervised pre-training, we use the full 960 hours of LibriSpeech audio~\cite{panayotov2015librispeech} or 60,000 hours of Libri-light~\cite{kahn2020libri} audio, both of which are derived from the LibriVox project that contains English recordings of copyright-free audiobooks by volunteers from the Internet.
For supervised fine-tuning, five different partitions are considered: Libri-light 10-minute, 1-hour, 10-hour splits and LibriSpeech 100-hour (\texttt{train-clean-100}) and 960-hour (\texttt{train-clean-100}, \texttt{train-clean-360}, \texttt{train-other-500} combined) splits. The three Libri-light splits are subsets of the the LibriSpeech training split, and each of them contain half of the audio from \texttt{train-clean-*} and the other from \texttt{train-other-500}.


\subsection{Unsupervised Unit Discovery}
To demonstrate the effectiveness of the proposed method on utilizing low-quality cluster assignments, we consider the k-means algorithm~\cite{lloyd1982least} for acoustic unit discovery by default. It is one of the most naive unit discovery models that can be treated as modeling an isotropic Gaussian with the same scalar variance for each acoustic unit.
To generate labels for the first iteration HuBERT training over the 960 hour LibriSpeech training set, we run k-means clustering with 100 clusters on 39-dimensional MFCC features, which are 13 coefficients with the first and the second-order derivatives.

To generate better targets for the subsequent iterations, we run k-means clustering with 500 clusters on the latent features extracted from the HuBERT model pre-trained in the previous iteration (not fine-tuned) at some intermediate transformer layer.
Since the feature dimension at the transformer output is much higher than the MFCC features (768-D for HuBERT \textsc{Base}), we cannot afford to load the entire 960 hour training split to the memory. So instead, we randomly sample 10\% of the data for fitting the k-means model.

The \texttt{MiniBatchKMeans} algorithm implemented in the \texttt{scikit-learn} \cite{pedregosa2011scikit} package is used for clustering, which fits a mini-batch of samples at a time.\footnote{It still requires loading the entire dataset to the memory first.} We set the mini-batch size to be 10,000 frames. k-means++~\cite{arthur2006k} with 20 random starts is used for better initialization.

\subsection{Pre-Training}\label{sec:pretrain}
We train the \textsc{Base} model for two iterations on the 960 hours of LibriSpeech audio on 32 GPUs, with a batch size of at most 87.5 seconds of audio per GPU. The first iteration is trained for 250k steps, while the second iteration is trained for 400k steps using labels generated by clustering the 6-th transformer layer output of the first iteration model. Training for 100k steps takes about 9.5 hours.

Next we train HuBERT \textsc{Large} and \textsc{X-Large} for one iteration on 60,000 hours of Libri-light audio on 128 and 256 GPUs, respectively, for 400k steps. The batch sizes are reduced to 56.25 and 22.5 seconds of audio per GPU due to memory constraints.
Instead of restarting the iterative process from clustering MFCC features, we extract features from the 9-th transformer layer of the second iteration \textsc{Base} HuBERT for clustering and use those labels for training these two models. Hence, these two models can also be seen as the third iteration models.

For all HuBERT configurations, mask span is set to $l=10$, and $p=8\%$ of the waveform encoder output frames are randomly selected as mask start if not otherwise mentioned. Adam~\cite{kingma2014adam} optimizer is used with $\beta = (0.9, 0.98)$, and the learning rate ramps up linearly from 0 to the peak learning rate for the first 8\% of the training steps, and then decays linearly back to zero. The peak learning rates are 5e-4/1.5e-3/3e-3 for \textsc{Base}/\textsc{Large}/\textsc{X-Large} models.


\subsection{Supervised Fine-Tuning and Decoding}
We fine-tune each model on 8 GPUs on the labeled splits described in Section~\ref{sec:data}. The batch sizes per GPU are at most 200/80/40 seconds of audio for \textsc{Base}/\textsc{Large}/\textsc{X-Large} models. During fine-tuning, the convolutional waveform audio encoder parameters are fixed. Like wav2vec 2.0, we introduce a \textit{freeze-step} hyperparameter to control how many fine-tuning steps the transformer parameters are fixed, and only the new softmax matrix is trained. 
We sweep over peak learning rate ([1e-5, 1e-4]), learning rate schedule (percentage of steps for linear ramp-up and decay), number of fine-tuning steps, freeze step, and waveform encoder output masking probability for each model size and fine-tuning split combination using the word error rate (WER) on the \texttt{dev-other} subset as a criterion for model selection.

We use the wav2letter++~\cite{pratap2018wav2letter++} beam search decoder wrapped in Fairseq~\cite{ott2019fairseq} for language model-fused decoding, which optimizes:
\begin{equation}
    \log p_{CTC}(Y \mid X) + w_1 \log P_{LM}(Y) + w_2 |Y|,
\end{equation}
where $Y$ is the predicted text, $|Y|$ is the length of the text, and $w_1$ and $w_2$ denote the language model weight and word score. The decoding hyperparameters are searched with Ax, a Bayesian optimization toolkit,\footnote{\url{https://github.com/facebook/Ax}}. In this work, we consider both $n$-gram and transformer language models trained on the official Librispeech language modeling data.


\subsection{Metrics of Target Quality}
For analysis, we derive frame-level forced-aligned phonetic transcripts using a hybrid ASR system to measure the correlation between the k-means cluster assignments and the actual phonetic units. 
Given aligned frame-level phonetic labels $[y_1, \cdots, y_T]$ and k-means labels $[z_1, \cdots, z_T]$, the joint distribution between the two variables $p_{yz}(i, j)$ can be estimated by counting the occurrences: 
\begin{equation}
    p_{yz}(i, j) = \dfrac{\sum_{t=1}^T [y_t = i \wedge z_t = j] }{T},
\end{equation}
where $i$ denotes the $i$-th phoneme class and $j$ denotes the $j$-th k-means label class.
The marginal probabilities are computed as $p_z(j) = \sum_i p_{yz}(i, j)$ and $p_y(j) = \sum_j p_{yz}(i, j)$.

For each phone class $i$, we further compute the most likely target label as:
\begin{equation}
    z^*(i) = \arg\max_j p_{yz}(i, j).
\end{equation}
Likewise, for each k-means class $j$, we compute the most likely phone label as:
\begin{equation}
    y^*(j) = \arg\max_i p_{yz}(i, j).
\end{equation}
Three metrics are considered:
\begin{enumerate}
    \item \textbf{phone purity} (Phn Pur.): 
    \begin{equation}
        \mathbb{E}_{p_z(j)} [ p_{y \mid z}(y^*(j) \mid j) ],
    \end{equation}
    where $p_{y \mid z}(i \mid j) = p_{yz}(i, j) / p_z(j)$ denotes the conditional probability of phone given a k-means label. This metric measures the average phone purity within one class, which can be interpreted as the frame-level phone accuracy if we transcribe each k-means class with its most likely phone label. When comparing different sets of target labels with the same number of units, higher purity indicates better quality. However, this metric is less meaningful when comparing two sets with different numbers of units: in the extreme case where each frame is assigned a unique target label, the phone purity would be 100\%.
    %
    \item \textbf{cluster purity} (Cls Pur.):
    \begin{equation}
        \mathbb{E}_{p_y(i)} [ p_{z \mid y}(z^*(i) \mid i) ],
    \end{equation}
    where $p_{z \mid y}(j \mid i) = p_{yz}(i, j) / p_y(i)$ denotes the conditional probability of a k-means label given phone label. Cluster purity is the counterpart of phone purity, whose value would typically decrease when the number of units increases. When comparing target labels with the same number of units, higher cluster purity also indicates a better quality, as frames of the same phone are more likely labeled as the same k-means label class.
    %
    \item \textbf{phone-normalized mutual information} (PNMI):
    \begin{align}
        \dfrac{I(y; z)}{H(y)} &= \dfrac{ 
            \sum_i \sum_j p_{yz}(i, j) \log \dfrac{p_{yz}(i, j)}{p_{y}(i)p_{z}(j)}
        }{
            \sum_i p_{y}(i) \log p_{y}(i)
        } \\
        &= \dfrac{H(y) - H(y \mid z)} {H(y)} \\
        &= 1 - \dfrac{H(y \mid z)} {H(y)}.
    \end{align}
    PNMI is an information-theoretic metric that measures the percentage of uncertainty about the phone label $y$ eliminated after observing the k-means label $z$. Higher PNMI also indicates better k-means clustering quality.
\end{enumerate}
